\section{DNS}
\label{chapter-DNS}
Im Internet kennen sich Computer über IP-Adressen.
Da es für Menschen nicht realistisch ist, sich für jeden Computer der angesprochen werden soll, sich die IP-Adresse des Computers zu merken,
gibt es das Domain Name System (DNS).
DNS ist ein ein System, das dazu da ist, IP-Adressen in Namen umzuwandeln die für Menschbar lesbar und umgänglich sind.
Ebenso kann DNS auch Namen in IP-Adressen umwandeln, was als reverse-DNS bezeichnet wird.
An einem praktischen Beispiel lässt sich das einfach zeigen.
Z.B. steht die IP-Adresse von example.com mit \textbf{host example.com} abrufen, in Abbildung-\ref{fig:output-host-example.com} ist ein DNS-Eintrag den man mit \enquote{host} (unter linux) abrufen kann, alternativ lässt sich dies auch mit z.B. \enquote{dig} auflösen.
\begin{figure}[htbp]
  \centering
  \lstinputlisting{src/host_example-com.txt}
  \caption{output von \enquote{host example.com}}
  \label{fig:output-host-example.com}
\end{figure}

DNS ist ein wichtiger Dienst um kommunikation im Internet zu ermöglichen, meistens ist es sogar einer der ersten Dienste die überhaupt die Kommunikation im Internet unter Computern ermöglicht.\\

DNS wurde entworfen möglichst leichtgewichtig zu sein, mit kleinen Paketen und wenig overhead (durch UDP), und dadurch auch geringeren Latenzen.
Das DNS Protokoll läuft standartmäßig auf UDP mit Port 53. TPC-fallback ist möglich, wenn die Nachrichten wegen der Größe fehlschlagen sollten, es ist aber auch möglich direkt DNS Pakete über TCP zu senden, mit gesetztem Bit\cite{ietf-rfc-1034,ServiceNameTransport}.
Die UPD Nachrichtengröße ist auf 512 Byte beschränkt, für DNS Nachrichten reicht das eigentlich aus\cite{ietf-rfc-1034,ietf-rfc-1035}.
Die Domänen werden wie in \ref{fig:dns-baum} zu sehen ist, in einer Baumstruktur angeordnet.

% diese figure ist von claude generiert worden. Muss erwähnt werden!
\begin{figure}[htbp]
\centering
\begin{forest}
for tree={
  rectangle, draw, minimum height=1.6em, minimum width=2.5em,
  l sep=25pt, s sep=15pt
}
[. (Root)
  [com
    [example.com
      [www.example.com]
      [mail.example.com]
    ]
  ]
  [de
    [example.de]
  ]
  [org
    [example.org]
  ]
]
\end{forest}
\caption{Beispiel der hierarchischen Struktur des DNS-Namensraums}
\label{fig:dns-baum}
\end{figure}

Der DNS Baum teilt sich anhand dessen Domänen auf. Die oberste Domäne ist die sogenannte root-Domain und wird mit einem \enquote{.} gekennzeichnet.
Unter der root-Domain liegen die Top-Level-Domains (TLD).
In Abbildung-\ref{fig:dns-baum} sind com, de und org die TLDs, es sind immer die Domänen, die unmittelbar unter der root-Domain liegen.
Die Knoten darunter, die keine Blätter sind, sind weitere Domänen, wenn sie weitere Kinder haben, wie in Abbildung-\ref{fig:dns-baum}, com, de, org und auch example.com, dann werden diese als Subdomänen bezeichnet.
Die Blätter werden als \enquote{labels} bezeichnet.
Wie die DNS-Server verteilt sind ist unterschiedlich, ihre zuordnung ist rein logisch, Domänen-Server sind verteilte Datenbanksysteme und die einzelnen Computer können unterschiedlich auf der ganzen Welt verteilt sein\cite{baunComputernetzeKompaktPraxisorientierte2025}.

In Abbildung-\ref{fig:dns-baum} ist eine mögliche Struktur einer hierarchischen Struktur eines DNS-Namensraumes zu sehen.
Die Namensauflösung geschieht von rechts nach links.
Die jeweilige Zone kann mehrere labels haben oder wiederrum subdomänen.
Im Beispiel von Abbildung-\ref{fig:dns-baum} haben org und de, jeweils ein label namens example.
Das der Name example in beiden Domänen vorkommt, führt zu keinem Konflikt.
Trotz dass der Name der gleiche ist, kann unterschieden werden, weil dieser nicht in der selben Domäne ist.
Im Fall von der com Domäne ist example wieder eine subdomäne, die die labels \enquote{www} und \enquote{mail} besitzt.